\begin{figure}[tbp]
\centering
\begin{tikzpicture}[
  font=\small,
  box/.style={draw=black!65,rounded corners=2pt,align=center,
    inner sep=6pt,text width=3.1cm,minimum height=1.05cm},
  flow/.style={-{Latex[length=2mm]},thick}
]
\node[box,fill=black!3] (input) at (0,0)
  {Canonical $(n,k,d)$\\Request counter $t$};
\node[box,fill=blue!5] (simulate) at (4.7,0)
  {Capped simulation\\of constructor $P$\\$g(t)$ transitions};
\node[box,fill=orange!8] (failure) at (9.4,0)
  {Serve with label $1$};
\node[box,fill=blue!5] (output) at (4.7,-2.4)
  {Completed table\\and parameters\\$S=O(g(t)+1)$};
\node[box,fill=blue!5] (suffix) at (9.4,-2.4)
  {First success:\\activate\\Later successes:\\advance state};
\node[align=center,text width=3.7cm,font=\footnotesize] (state) at (9.4,-4.05)
  {Retain suffix state\\for the next request.};
\node[align=left,text width=3.6cm] at (0,-2.4)
  {$g(t)=\lfloor\log_2(t+1)\rfloor$\smallskip\\
   First success:\\$t_*=2^{T_P}-1$.};
\draw[flow] (input) -- (simulate);
\draw[flow] (simulate) -- node[above,align=center,font=\footnotesize]{not\\halted} (failure);
\draw[flow] (simulate) -- node[right,font=\footnotesize]{halted} (output);
\draw[flow] (output) -- (suffix);
\draw[flow] (suffix.south) -- (state.north);
\end{tikzpicture}
\caption{Each request pays for a fresh capped simulation. Before activation,
only one fixed label serves requests. After activation, the complete table's
explicit size fits the current budget, and the suffix routine advances its
persistent state, with resets only at its prescribed epoch boundaries.
The construction time $T_P$ and activation time $t_*$ are
finite instance-dependent quantities; their magnitude enters the additive
movement term.}
\label{fig:activation}
\end{figure}
